% Pregunta 3. Consumo, ahorro, inversión y crédito (bloque II.2) · peso medio
\section*{Pregunta 3. Consumo, ahorro, inversión y crédito}

\begin{enunciado}
\textbf{3.} Analice el comportamiento del \textbf{consumo, el ahorro y la inversión} del sector privado, así como sus principales determinantes. ¿Qué factores podrían explicar el dinamismo de la inversión en los últimos años y su distinta composición respecto de la fase expansiva anterior a 2008? Valore la evolución del \textbf{endeudamiento y del crédito} al sector privado.

\smallskip
\textbf{Peso medio} · \tiempo{17} · guía: I.4 (consumo e inversión: determinantes), II.2 (crédito), IV.1 (ahorro e inversión), V.2 (endeudamiento, vivienda) · cuadros 1 y 4; gráficos 4 y 5.
\end{enunciado}

\begin{examen}
\textit{El sector privado español ha pasado de gastar a crédito, como en 2007, a ahorrar más de lo que invierte: los hogares ahorran el doble que antes de la crisis, la inversión crece con otra composición (intangibles e infraestructuras, no vivienda) y la deuda privada se ha reducido a la mitad. El único foco de riesgo es el precio de la vivienda.}

\rotulo{a) Consumo, ahorro e inversión y sus determinantes}
La tasa de ahorro de los hogares (cuadro~4, gráfico~5) resume los episodios del ciclo: sube del 6,0 al 11,7~\% entre 2007 y 2009 por \emph{precaución}; baja al 7,0~\% en 2012, cuando la renta disponible real por habitante cae un 6,6~\% y los hogares desahorran para sostener el consumo; se dispara al 17,6~\% en 2020 por el ahorro \emph{forzoso} y el precautorio; y se gasta en 2021-2022 (9,0~\%), cuando la inflación erosiona la renta real ($-2{,}0~\%$ en 2022). Desde 2023 la renta real se recupera (+3,2 y +3,3~\% por habitante en 2023-2024) más deprisa que el consumo, y el ahorro sube al 12,8~\%. En 2025 ocurre lo contrario: el consumo crece un 3,1~\% con una renta por habitante que solo avanza un 0,9~\%, y la tasa de ahorro baja al 11,9~\%. Entre los determinantes habituales, el consumo se apoya hoy en el \textbf{empleo} (+2,6~\%) y en la \textbf{población}, en la riqueza inmobiliaria y en unos tipos que bajaron en 2025; frena la renta real por persona. Que el ahorro doble el de 2007 sería coherente con el desapalancamiento, la incertidumbre y el envejecimiento.

La tasa de inversión de las empresas (24,4~\% del valor añadido) está lejos del 30,2~\% de 2007, y la de los hogares en vivienda (7,4~\% de su renta), muy lejos del 15,8~\%. Los determinantes son mixtos: la \textbf{brecha positiva} y la demanda esperada animan a ampliar capacidad (acelerador); el \textbf{coste de la financiación} subió con fuerza en 2023-2024 (préstamos nuevos a empresas al 4,6-4,7~\%) y baja en 2025 (3,4~\%, apenas 0,4 puntos por encima del deflactor del PIB); la \textbf{rentabilidad} empresarial se estrecha (los beneficios pasan del 43,8~\% del valor añadido en 2016 al 38,9~\%); y la \textbf{incertidumbre} (energía, aranceles) juega en contra.

\rotulo{b) El dinamismo de la inversión y su composición}
La formación bruta de capital fijo crece un 13,4~\% entre 2019 y 2025, pero con una composición muy distinta de la de antes de 2008: los \textbf{productos de la propiedad intelectual} (programas informáticos, I+D, bases de datos) aumentan un 47,4~\% y las \textbf{otras construcciones} (infraestructuras) un 14,9~\%, mientras la vivienda solo crece un 6,0~\% y la maquinaria un 2,7~\% (cuadro~1). Antes de 2008 tiraba la vivienda, que se hundió un 51,6~\% entre 2007 y 2013. Detrás del cambio habría 3 factores: la \textbf{digitalización}, que desplaza la inversión hacia los intangibles; los \textbf{fondos europeos}, que llegan como transferencias de capital (la cuenta de capital ronda el 1~\% del PIB desde 2021, cuadro~7) y financian infraestructuras, digitalización y transición energética; y una \textbf{oferta de vivienda rígida} (suelo, plazos, menos crédito promotor), que explica que la inversión residencial apenas crezca pese a la demanda. Es una inversión más sana que la de 2007, financiada con ahorro nacional (la tasa de ahorro nacional sube al 25,0~\% del PIB y supera a la inversión, del 21,7~\%) y que eleva el potencial, porque \uline{la inversión es demanda hoy y oferta mañana}.

\rotulo{c) Endeudamiento y crédito}
El \textbf{desapalancamiento} ha sido profundo y ordenado. El flujo de crédito al sector privado pasa del 26,5~\% del PIB en 2007 al $-11{,}2$~\% en 2012 y al 2,6~\% en 2025 (cuadro~4); la deuda de los hogares baja del 85,0~\% del PIB en 2009 al 42,9~\%, y la de las empresas, del 119,0 al 62,0~\% (gráfico~4), ambas por debajo de los umbrales del procedimiento de desequilibrios (55 y 85~\%). Desde 2019 el ajuste ha sido sobre todo por el denominador: la deuda de los hogares apenas varía en euros y su ratio baja porque el PIB nominal crece un 35~\%. Es un \emph{desapalancamiento sano}, hecho con más renta y no hundiendo el gasto. El crédito, que cayó en 2023 ($-2{,}0~\%$ a hogares y $-3{,}7~\%$ a empresas) con la subida de tipos, se reactiva en 2025 (+4,5 y +4,2~\%): la política monetaria se transmite y no hay racionamiento.

La excepción es la \textbf{vivienda}: su precio sube un 12,7~\% en 2025 (9,7~\% en términos reales), por encima del umbral del 9~\% del procedimiento de desequilibrios, aunque en términos reales sigue un 13~\% por debajo de 2007 (gráfico~4). A diferencia de entonces, no la impulsa el crédito, sino una demanda que crece con la población frente a una oferta que no responde: es más un problema de \emph{accesibilidad} que de estabilidad financiera, aunque conviene vigilarlo (pregunta~8).
\end{examen}

\begin{complemento}{3.1 · Por qué se contesta así: los determinantes, con su signo y su dato}
\que{Qué se aplica.} Guía, I.4: las funciones de consumo e inversión con sus signos, eligiendo los 2 o 3 determinantes que el cuadro permite ver y diciendo cuáles faltan (nota al opositor de I.4); II.2 (crédito como factor); IV.1 ($S-I$); V.2 (desapalancamiento sano, vivienda).

\begin{tablacaja}
\centering\scriptsize
\caption{Los determinantes del consumo y la inversión en el cuadernillo (2023-2025)}\label{tab:determinantes}
\begin{tabular}{>{\raggedright\arraybackslash}p{3.1cm}>{\centering\arraybackslash}p{0.9cm}>{\raggedright\arraybackslash}p{6.0cm}>{\raggedright\arraybackslash}p{3.6cm}}
\toprule
\rowcolor{tceeazulpalido}
\textbf{Determinante} & \textbf{Signo} & \textbf{Dato} & \textbf{Efecto} \\
\midrule
\rowcolor{tceesalvia}\multicolumn{4}{l}{\emph{Consumo}: $C=f(Y^d_{\text{real}}{}^{+},\ \text{riqueza}^{+},\ r^{-},\ \pi^{-},\ \text{crédito}^{+},\ \text{deuda}^{-},\ N^{+},\ \text{incertidumbre}^{-})$} \\
Renta disponible real & $+$ & +3,2, +3,3 y +0,9~\% por habitante (C4) & Sostiene 2023-2024; frena en 2025 \\
Empleo y población & $+$ & Empleo +2,6~\%, población +1,0~\% en 2025 (C3) & Tira del consumo total \\
Riqueza (vivienda) & $+$ & Precio +12,7~\% en 2025 (C4) & Empuja (efecto riqueza) \\
Tipo real & $-$ & Hipotecas al 2,73~\% con IPCA del 2,7~\%: tipo real $\approx0$ (C4, C5) & Neutro o favorable \\
Deuda de los hogares & $-$ & 42,9~\% del PIB, mínimo de la serie (C4) & Ya no frena \\
\rowcolor{tceesalvia}\multicolumn{4}{l}{\emph{Inversión}: $I=f(r^{lp\,-},\ \gamma^e_Y{}^{+},\ \mathrm{OG}^{+},\ \mathrm{EBE}^{+},\ \text{crédito}^{+},\ \text{incertidumbre}^{-})$} \\
Brecha y demanda & $+$ & Brecha +2,0~\%; PIB +2,6~\% (C1) & Empuja (acelerador) \\
Coste de financiación & $-$ & Empresas: 4,64~\% (2023) → 3,40~\% (2025); real $\approx3{,}40-3{,}0=0{,}4$ & Frena en 2023-2024, alivia en 2025 \\
Rentabilidad & $+$ & Beneficios 40,3~\% del VAB (2022) → 38,9~\% (C4); margen $-1{,}4$ en 2025 & Frena \\
Fondos europeos & $+$ & Cuenta de capital 0,9-1,1~\% del PIB desde 2021 (C7) & Empuja (infraestructuras, intangibles) \\
Incertidumbre & $-$ & No hay indicador; energía y aranceles (G7, G9) & Se dice con condicional \\
\bottomrule
\end{tabular}
\fuente{Cuadernillo del examen (C: cuadro; G: gráfico). Signos: guía, I.4. El tipo real a largo se aproxima con el tipo de los préstamos nuevos menos el deflactor del PIB (la guía, I.4, nota, da el orden de preferencia).}
\end{tablacaja}

\que{Error típico.} Explicar la caída del ahorro de 2025 con «los hogares consumen más porque tienen más renta»: es al revés, el consumo crece \emph{más} que la renta por habitante; lo explican el empleo, la población y la riqueza. Y decir que el consumo crece «un 3,1~\% por habitante»: el 3,1 es total; por habitante, en torno al 2~\% ($3{,}1-1{,}0$).
\end{complemento}

\begin{complemento}{3.2 · Los cálculos: inversión por activos, ahorro e inversión y deuda}
\que{Inversión acumulada por activos} (encadenando las tasas del cuadro~1):

\begin{tablacaja}
\centering\scriptsize
\caption{Formación bruta de capital fijo por activos: variación acumulada en volumen (\%)}\label{tab:fbcf}
\begin{tabular}{>{\raggedright\arraybackslash}p{4.6cm}*{4}{>{\centering\arraybackslash}p{1.8cm}}}
\toprule
\rowcolor{tceeazulpalido}
\textbf{Activo} & \textbf{2007-2013} & \textbf{2013-2019} & \textbf{2019-2025} & \textbf{2007-2025} \\
\midrule
Formación bruta de capital fijo & $-37{,}4$ & 33,8 & 13,4 & $-5{,}0$ \\
\hspace*{0.8em}Viviendas & $-51{,}6$ & 61,9 & 6,0 & $-16{,}9$ \\
\hspace*{0.8em}Otras construcciones & $-44{,}6$ & 14,8 & 14,9 & $-27{,}0$ \\
\hspace*{0.8em}Maquinaria, bienes de equipo y armamento & $-28{,}7$ & 36,3 & 2,7 & $-0{,}2$ \\
\hspace*{0.8em}Productos de la propiedad intelectual & 19,7 & 21,7 & 47,4 & 114,7 \\
PIB & $-7{,}6$ & 16,9 & 10,2 & 19,0 \\
\bottomrule
\end{tabular}
\fuente{Cálculo propio con el cuadro~1 ($\prod_t(1+\gamma_t)-1$; \texttt{calculos.txt}). Qué hay que ver: la inversión total aún no ha recuperado el nivel de 2007; solo los intangibles lo superan (más que lo duplican).}
\end{tablacaja}

\que{Ahorro, inversión y saldo exterior} (guía, IV.1 y A.3.4): $S/\mathrm{PIB}=I/\mathrm{PIB}+\mathrm{CC}/\mathrm{PIB}$. En 2025, $25{,}0-21{,}7=3{,}3\approx3{,}2$ (cuenta corriente, cuadro~7); en 2007, $21{,}1-30{,}5=-9{,}4$, exactamente el déficit corriente. El paso de un $-9{,}4$ a un +3,2 vino primero de la inversión (cae 13,1 puntos entre 2007 y 2013) y después del ahorro (sube 5,5 puntos entre 2013 y 2025): enlaza con la pregunta~6.

\que{Desapalancamiento: numerador o denominador.} Deuda en euros $=$ ratio $\times$ PIB nominal (cuadro~1). Hogares: $0{,}565\times1.254\approx708$ miles de millones en 2019 y $0{,}429\times1.690\approx725$ en 2025 (+2~\%); si la deuda no hubiera variado en euros, la ratio habría bajado sola al 41,9~\%. Entre 2009 y 2019, en cambio, la deuda en euros cayó un 22~\%: \uline{primero un desapalancamiento activo (amortizar), después uno pasivo (crecer)}. Empresas: +10,6~\% en euros entre 2019 y 2025 (los préstamos con aval público de 2020: préstamos a empresas +9,1~\% ese año).

\que{Vivienda} (cuadro~4): precio nominal acumulado 2013-2025, +88~\%; real, +48~\%; real 2019-2025, +19~\% con una inversión en vivienda que solo crece un 6,0~\%: \emph{más precio que cantidad}, la señal que la guía (V.2) asocia a tensiones.

\que{Qué añadir con más tiempo.} La relación entre precio de la vivienda y renta (no viene en el cuadernillo) y el esfuerzo de acceso (umbral orientativo de un tercio de la renta, guía V.2). \que{Qué recortar en el examen.} Las tasas de inversión de hogares y empresas (basta la composición por activos) y la descomposición del desapalancamiento (basta «la ratio baja sobre todo porque crece el PIB nominal»).
\end{complemento}
